\section{The harmonic equality barrier and the uniform gap}\label{sec:equality}

The moment identity and strict transport can coexist at
\(\beta\ge1/4\) only in the equality case \(\beta=1/4\).
We first rule out its deterministic consequence, an identically
vanishing Jacobian, and then transfer that obstruction back to the law.

\begin{lemma}[The harmonic equality barrier]\label{lem:no-zero-J}
For \(k=3/4\), no \(F\in\KK\) has \(J_F\equiv0\).
\end{lemma}
\begin{proof}
By \eqref{eq:Jacobian}, such a map would satisfy
\[
 |q_F(z)+1/4|^2=1/4.
\]
Therefore \(|q_F(z)|\ge1/4\) and
\(|F''(z)/F'(z)|\ge1/(4y)\).
The pre-Schwarzian has no zeros, so
\(u=\log|F''/F'|\) is harmonic and \(u\ge-\log4-\log y\).
For \(0<\epsilon<1/2\), the harmonic function
\[
 b_\epsilon(x,y)=-\log4+
 \log(1/\epsilon)\cos(\pi x/2)
 \frac{\sinh(\pi(1-y)/2)}{\sinh(\pi/2)}
\]
lies below \(u\) on the boundary of
\([-1,1]\times[\epsilon,1]\).
On its sides and top the added term is zero; on its bottom
the cosine-sinh factor is between zero and one.
The minimum principle gives \(u\ge b_\epsilon\) inside.
At \(i/2\), the coefficient of \(\log(1/\epsilon)\) is fixed
and positive, contradicting finiteness as \(\epsilon\downarrow0\).
\end{proof}

\begin{proof}[Proof of Theorem~\ref{thm:main}]
Under \eqref{eq:contrary}, Equation~\eqref{eq:Jacobian} and
Proposition~\ref{prop:strict-transport} force
\(\beta=1/4\), \(k=3/4\), and \(J_F(i)=0\) almost surely.
Apply \eqref{eq:law} to the test \(|J_F(i)|\).
By \eqref{eq:G-action} and the strictly positive weight
\(|Q_F(z)|\), \(J_F(z)=0\) almost surely at every fixed \(z\).
A countable dense set and continuity give \(J_F\equiv0\)
for almost every \(F\), contradicting Lemma~\ref{lem:no-zero-J}.
Thus the single exponent in \eqref{eq:rho} satisfies \(\beta<1/4\).

Choose \(b\) with \(\beta<b<1/4\).
The root limit in \eqref{eq:rho} supplies a constant \(C_b\)
such that \(\|L^n\|\le C_b2^{bn}\) for all \(n\).
Proposition~\ref{prop:disk}, with \(2^{-n}\le1-r<2^{1-n}\),
then gives \(M_{-1}[f'](r)\le C(1-r)^{-b}\), uniformly in
\(f\in\SSS\). Taking \(\varepsilon=1/4-b\) proves
\eqref{eq:uniform} and the stated disproof.
\end{proof}

